\cotefile{29}
\scanpage{49}
% archive foliation: 46 (pencil, lower left; uncertain)
\begin{center}- 23 -\end{center}
\begin{flushright}n°385\end{flushright}

\uline{l'injection de} $Y$ \uline{dans} $X$. \uline{Pour tout point} $a$ \uline{de} $Y$, \uline{l'application} $\iota_a$ \uline{est une bijection de} $D_a(Y)$ \uline{sur le sous-espace de} $D_a(X)$ \uline{constitué des} $\mu$ \uline{telles que} $f.\mu=0$ \uline{pour toute fonction de classe} $C^r$ \uline{dans un voisinage} $U$ \uline{de} $a$, \uline{nulle en tout point de} $Y\cap U$.

Soit $\mu$ dans $D_a(Y)$ et $f$ de classe $C^r$ dans un voisinage $U$ de $a$, nulle sur $Y\cap U$ ; on a :
\[
f.\iota_a(\mu)=\iota_a((f\circ\iota).\mu)=0
\]
puisque $f\circ\iota=0$. Comme $Y$ est une sous-variété de $X$, il existe un ouvert $U$ contenant $a$, un système de coordonnées $x_1,\ldots,x_n$ dans $U$, et un entier $m$ tel que $0\leq m\leq n$ et que $Y\cap U$ soit l'ensemble des points de $U$ où s'annulent $x_{m+1},x_{m+2},\ldots,x_n$. Les fonctions $x_1\circ\iota,\ldots,x_m\circ\iota$ forment alors un système de coordonnées sur l'ouvert $Y\cap U$ de $Y$.

L'espace $D_a(Y)$ a une base $\{\Lambda^\alpha\}_{\alpha\in\mathbf{N}_r^m}$ caractérisée par :
\[
\langle\Lambda^\alpha,x^{\alpha'}\circ\iota\rangle=\delta_{\alpha\alpha'}
\qquad \alpha,\alpha'\in\mathbf{N}_r^m
\]
et l'espace $D_a(X)$ admet une base $\{\Delta^\beta\}_{\beta\in\mathbf{N}_r^n}$ définie par :
\[
\langle\Delta^\beta,x^{\beta'}\rangle=\delta_{\beta\beta'}
\qquad \beta,\beta'\in\mathbf{N}_r^n.
\]
Nous noterons $0$ l'élément nul de $\mathbf{N}_r^{n-m}$ ; on a alors $\iota_a(\Lambda^\alpha)=\Delta^{(\alpha,0)}$ pour tout $\alpha$ dans $\mathbf{N}_r^m$ et par conséquent, $\iota_a$ est injective.

Par ailleurs, soit $\mu$ dans $D_a(X)$ telle que $x_i.\mu=0$ pour $m<i\leq n$. D'après (59), on en déduit $\langle\mu,x^\beta\rangle=0$ pour tout multiindice $\beta$ dans $\mathbf{N}_r^n$ pour lequel il existe $i$ avec $m<i\leq n$ et $\beta_i\neq0$ ; autrement dit, on ne peut avoir $\langle\mu,x^\beta\rangle\neq0$ que si $\beta=(\alpha,0)$ pour $\alpha$ dans $\mathbf{N}_r^m$ convenable. De la formule $\mu=\sum_\beta\langle\mu,x^\beta\rangle.\Delta^\beta$, on déduit alors $\mu=\sum_\alpha\langle\mu,x^\alpha\rangle\,\iota_a(\Lambda^\alpha)$, ce qui prouve que $\mu$ est dans l'image de $\iota_a$.
\begin{flushright}C.Q.F.D.\end{flushright}

Pour tout point $a$ de $X$, l'espace vectoriel $D_a(X)\otimes D_a(X)$ est muni d'une structure
